\section{Discussion}
\label{sec:discussion}

\subsection{What the combined evidence establishes}

The implemented contribution is the coupling between a software progress boundary and all persistent representations of the retained conversational prefix. Transcript-only truncation can leave discarded content accessible through K/V state; tensor-only cropping can leave masks, positions, role structure, or ledgers at an incompatible endpoint. The two zero-content transitions and the pending-EOT lifecycle make these distinctions operational rather than merely descriptive. Joint consistency is checked at completed operation boundaries, not asserted as transactional atomicity for concurrent readers.

The evidence addresses complementary failure modes. C2 checks that the requested prefix survives cropping and matched recovery. B checks software interval selection against an independently constructed reference. R asks whether the repaired state can actually accept the same next-turn input and produce a deliverable token at lower operation cost than full rebuilding. The large long-context difference is consistent with avoiding retained-prefix recomputation. It is not a new asymptotic cache result, but an empirical link between correct conversational serialization and useful recovery work avoided in this implementation.

Structural validity must remain separate from numerical and response identity. R used different forward topologies for reuse and rebuilding: next-token choices agreed in every pair, yet all repetitions of one second-interruption cell produced different short continuations. This is observable output variation, not a tolerance-based equivalence pass. Likewise, the rejected C2 clean-reprefill protocols remain rejected; the exact-only protocol answers a different, identifiable question. Neither the cost advantage nor the structural gates imply that arbitrary subsequent responses are interchangeable.

\subsection{Boundary policy and downstream uncertainty}

Whole-fragment retention deliberately favors keeping the consumed portion of an interrupted fragment at the cost of also retaining its unconsumed tail. Dropping the fragment would make the opposite error. The partial flag makes this tradeoff visible but does not solve it. B verifies application of the declared policy; it does not show that the policy is perceptually optimal or that simulated sample geometry corresponds to real speech alignment.

E3 is relevant precisely because a correct implementation need not improve dialogue. All four automated point estimates were negative in the generation-minus-cursor direction, but their intervals crossed zero. No substantive equivalence or non-inferiority margin was specified. The lexical and model-based measurements capture different forms of information reproduction and agree imperfectly; exact-key deduplication removes repeated weights, not semantically similar responses. Without blinded human reference labels, E3 cannot establish perceived coherence, trust, or task success. Its inconclusive result therefore constrains the contribution rather than being replaced by a favorable construction check.

\subsection{Validity and transfer conditions}

R's four process sessions and fixed cases provide limited replication, not a workload sample or an operational failure-rate estimate. Reconciliation before the second interruption holds the compared prefixes fixed but excludes free-running divergence across many turns. First-deliverable costs are sums of synchronized intervals, excluding preparation and CPU audit pauses. They should not be interpreted as uninterrupted response times or compared directly with P1's nested stop-path intervals. A1 fixes operation order and suffix length, while P1 has only 20 observations per cell; neither supplies a production tail-latency guarantee. The comparator in R reconstructs the full retained prefix and does not represent all serving systems with shared or stable-prefix caches.

The shared external limitation is the observed boundary. Software-consumed samples, device-presented samples, and acoustically heard content differ because of queues, APIs, scheduling, device buffers, propagation, and the listening environment. No reported experiment closes a real streaming-ASR, asynchronous-TTS, and audio-device loop or measures a loopback waveform. The evidence is concentrated on English task-oriented text, a frozen Qwen2-7B artifact, a ChatML-like template, Transformers DynamicCache, BF16/SDPA, and RTX 3090 hardware. Other templates, precisions, backends, languages, or segmentation rules require renewed boundary and state tests.

A transferable implementation therefore needs explicit fragment-to-chunk ownership, an interpretable playback-progress signal, and a backend that retains or reconstructs a legal prefix while updating mask, positions, ledgers, content span, and role/EOT state together. Word-level durations, device clocks, or loopback waveforms could refine the observed boundary, but would change what is being validated. Production claims additionally require concurrency and cancellation tests; interaction-quality claims require direct human evidence. These are extensions of the bounded method, not unmeasured benefits of the current implementation.

\section{Conclusion}
\label{sec:conclusion}

Joint prefix-state repair turns a software playback boundary into a checkable conversational-state transition. Its implementation distinguishes complete candidate invalidation from zero-content playback, preserves one consistent token/cache prefix, and commits structural closure once. Same-snapshot validation found all 27 crops and 60 matched recovery steps exact, while independent software replay checked 1118 cursor queries over 100 trajectories.

In the executable same-policy comparison, crop repair avoided full-rebuild work through the next-turn first-deliverable endpoint. At ordinary 512-, 2048-, and 8192-token boundaries the mean paired operation-cost differences were 113.621, 465.377, and 1956.100~ms. This is bounded recovery-cost evidence, not continuous speech end-to-end latency or a comparison against every prefix-caching strategy. Short continuations differed in 16/160 pairs despite identical first-token choices, and the downstream automated diagnostic remained inconclusive. The resulting contribution is an inspectable, incrementally useful systems method with explicit correctness and cost boundaries, rather than a new playback principle or a claim of numerical or perceptual equivalence.
